\chapter{The Linux Underneath}\label{cha:linux}

Every emulated machine sits on a real one. On most of them that fact is
carefully hidden; here it is a feature, because the real machine has
compilers on it and the fantasy one is being used to write programs.

This chapter is about the seam --- the one character that crosses it and
the words that use it --- and then about the Linux itself: what is on
it, how it gets onto a stick or beside another system, and how it boots.

\section{\texttt{!} --- the Linux prompt, from the machine's prompt}
A line beginning with \verb|!| is not the machine's:

\begin{type}
!ls -l
!git status
!nvim notes.txt
\end{type}

\noindent runs that command on the Linux the machine is standing on, with
its output drawn on the machine's screen. A bare \verb|!| with nothing
after it gives you an interactive shell there; leave it the way you leave
any shell.

It is carried on the Tube (Chapter~\ref{cha:tube}), which is the part of
this machine already built to have another processor's console on the
glass: the host's shell is started on the same pty the co-processors use,
JIM renders it, and the ROM's own key loop feeds it. So the colours stay
the machine's, full-screen programs work --- \verb|nvim|, \verb|htop| and
\verb|tmux| run in the window --- and \verb|ls --color| lands in the
machine's palette. The shell is told its terminal is
\textbf{xterm-color}, and it speaks UTF-8: JIM draws what it can in the
machine's character set and a near likeness of the rest, and an accented
letter you type arrives as one.

\verb|SSH [user@]host| is the same thing for a computer across the
network: an ssh session, drawn by JIM, through the Linux's own ssh.

\section{\texttt{NVIM} --- Neovim, dressed for the machine}
The machine's own editor is \verb|VI|. The Linux's is Neovim, and
\verb|NVIM [name]| brings it onto the glass the way \verb|!| would, but
set up for this machine (\pth{tools/nvim/init.lua}):

\begin{itemize}[leftmargin=1.4em]
\item It is in the machine's colours --- JIM's sixteen, which are the
  palette's, so a \verb|PALETTE| you load recolours it too --- and it
  draws nothing the code page has not got.
\item It knows the machine's languages by their extensions in capitals:
  \texttt{.C} is C (not the C++ Neovim takes \texttt{*.C} for),
  \texttt{.PAS} Pascal, \texttt{.RX} REXX, \texttt{.BAS} the machine's
  BASIC, \texttt{.BBC} BBC BASIC, \texttt{.LGO} LOGO. The BASIC's and
  LOGO's keywords are read from the interpreters' own sources
  (\pth{tools/mknvim.py}), so a keyword added there is coloured here.
\item \verb|:make| (or F9) compiles with the machine's compilers --- the
  same \verb|CC| and \verb|PAS|, a \texttt{PROJECT.K4P} beside the file
  included --- and reads \pth{/SYSTEM/LOG/MAKE.ERR}, the file
  \verb|VI|'s \verb|:make| reads, into Neovim's error list: it puts you on
  the first error, and \verb|:cn|, \verb|:cp| and \verb|:cl| go through
  them as in \verb|VI|.
\item \verb|:Run| (or F10) builds it and the machine runs it: Neovim
  steps aside, the program runs on the machine as \verb|VI|'s \verb|:run|
  runs it --- REXX, BASIC and LOGO in their interpreters --- and a key
  brings Neovim back, on the line you were on. A REXX program that stops
  on an error brings you back to that line.
\end{itemize}

\noindent Everything else is Neovim as it comes.

\begin{aside}
\textbf{Locking the door.} On a machine a child uses, the way into Linux
may be a way into trouble. \verb|linux = locked| in the menu file
(Section~\ref{sec:menufile}) turns \verb|!|, \verb|SSH| and the menu's
telnet row away with a message, and \verb|consoles = locked| makes the
Linux consoles unreachable from the keyboard. \verb|PAS| and \verb|CC|
still compile.
\end{aside}

\begin{aside}
\textbf{What this means for the container.} The sandboxed flavour
(\pth{linux/podman.sh}) is a wall on purpose: it is handed the display,
the sound, the game controllers and one folder, which the machine sees
as \pth{/MNT/SHARE}, and nothing else. \verb|!| inside it is real, and
reaches the container --- which is a small Debian with the compilers on
it, and not your computer.
\end{aside}

\section{\texttt{PAS} and \texttt{CC}: compiling on the machine}
The compilers live on the Linux side, so the machine reaches them the
same way you would --- but without you having to remember four options:

\begin{type}
VI HELLO.PAS
PAS HELLO
HELLO
\end{type}

\verb|PAS name| compiles \texttt{name.PAS} with Mad Pascal, and leaves
\texttt{name.prg} beside it; the name is looked for in the directory you
are standing in, or give a path, \verb|PAS /LANG/PASCAL/PMANDEL|.
\verb|CC name| does the same for \texttt{name.C} with cc65, using exactly
the rule the machine's own Makefile uses for its programs. A Pascal unit
beside the program is found whatever case its name was saved in (Mad
Pascal on Linux looks for \texttt{myunit.pas}; the machine writes
\texttt{MYUNIT.PAS}). Intermediates go to a scratch directory and never
appear on your disk. The compiler's errors appear on the screen, one to a
line as \texttt{FILE:LINE: Error: ...}, and in
\pth{/SYSTEM/LOG/MAKE.ERR}, which is what VI's \verb|:make| reads
(Chapter~\ref{cha:editors}); the exit status comes back as the result
code, so a script can test it --- as any \verb|!| command's does.

That is the loop the machine owns: \emph{edit the source on the machine,
compile it from the machine's prompt, run it on the machine.} The
compiler itself still runs on the Linux beside it, which is why
``self-hosted'' has an ``almost'' in front of it --- but the part that
matters when you are writing a program is all on this side of the seam.

\section{What is on the other side}
On a desktop, the other side is your own computer and has whatever you
have put on it. On the K4510's own Linux it is a Debian built for the
purpose, and the list is short because everything on it is there for a
reason:

\begin{center}
\small
\begin{tabular}{@{}l >{\raggedright\arraybackslash}p{0.58\linewidth}@{}}
\toprule
cc65, 64tass, nasm & the machine's own toolchain --- it can rebuild itself \\
FPC + Mad Pascal & what \verb|PAS| runs \\
git, neovim & so that work done here is work \\
curl, ssh & what \texttt{ftp://}, \texttt{sftp://} and \verb|SSH| go through \\
\texttt{tek40xx} & a Tektronix 4010 terminal, below \\
telnetd & bound to the machine itself and nothing else \\
NetworkManager, Tailscale & the network: F12 $\rightarrow$ Host \\
\bottomrule
\end{tabular}
\end{center}

\textbf{The consoles.} Ctrl+Alt+F1 is the machine; Ctrl+Alt+F2 to F6 are
Linux text consoles, 80 by 25 in the IBM PC's own font, where the
Tektronix lives; Ctrl+Alt+F1 comes back. (On a laptop whose F-keys are
media keys, add Fn.) Quitting the emulator (Shift+Esc, or F12
$\rightarrow$ Quit) drops you to a Linux shell on the machine's own
console rather than to nothing, and F12 has a \emph{Shut down the
computer} row because there is no desktop to go back to. The power
button performs a clean shutdown. Closing the lid does what F12
$\rightarrow$ Host $\rightarrow$ \emph{Lid closed} says: \emph{keep
running}, to begin with, or \emph{suspend}. (Quit to the Linux shell and
nothing is keeping the machine awake: there the lid suspends.) The consoles speak UTF-8 and follow
the keyboard layout chosen in F12.

\textbf{The network.} F12 $\rightarrow$ Host shows the computer's name
and address, and \emph{Wi-Fi / network setup} opens NetworkManager's
own screen on a spare console to join a network; the machine is back
when you leave it. A network joined once is remembered.

\textbf{Telnet, both directions.} \verb|TELNET 127.0.0.1 23| from the
machine's prompt --- or F12 $\rightarrow$ Host $\rightarrow$ \emph{Telnet
into the host} --- logs in to the Linux underneath through the front
door rather than through \verb|!|, which is useful when you want a
session that survives what the machine is doing. The daemon listens on
the machine itself only, so there is no network path to it at all.
Outbound, the network is real: \verb|TELNET| reaches a BBS, and URLs
work as Chapter~\ref{cha:shell} describes.

\section{On a stick}
\begin{type}
sudo ./linux/build-live.sh
\end{type}

\noindent makes the image, and \texttt{dd} puts it on a USB stick. At
boot a menu offers the keyboard layout, and then the whole system is
copied into RAM, so the stick is not being read while you work. A
persistence partition on the same stick keeps what should survive a
restart: the machine's home directory --- \texttt{k4510.cfg}, the menu
file, your saved states and the machine's disk with your files on it
--- the networks you have joined, and Tailscale's identity. Leave the
stick in.

\textbf{The laptop's own drive is not there.} The stick's kernel command
line blacklists the storage drivers absolutely, so the laptop's disk is
not mounted, not visible, and not writable by anything on the stick.
That is not a policy setting you could turn off by accident; the machine
cannot see it.

\section{Beside another system}\label{sec:install}
A computer that already runs Linux can have the machine as a second
choice at power-on. Make an empty partition of 2\,GB or more on its
internal disk and label it \texttt{K4510}; boot the computer's usual
Linux with the stick plugged in; and from the stick run

\begin{type}
sudo ./install-k4510.sh
\end{type}

\noindent It formats that partition the first time (and never again),
copies the system onto it, arranges for your settings and files to be
kept in the rest of the partition, and adds two entries to the
computer's boot menu:

\begin{description}[leftmargin=1.4em,style=nextline]
\item[K4510 Fantasy Computer] the machine, quietly: the machine's logo
  while Linux starts and nothing else, then the prompt.
\item[K4510 (text boot)] the same system, showing Linux's messages as
  it starts --- the one to choose when something needs looking at.
\end{description}

\noindent It refuses a partition on a USB disk, and one on the stick
itself, and it touches nothing of the other system but its boot menu.
Run it again from a newer stick to update: the system is replaced and
your settings and files are kept. The stick is only needed for the
install.

\textbf{The boot menu.} The computer boots whichever of the two systems
was used last, so the other system's own updates, which restart it,
come back to it. With the menu set to stay hidden, as it is on the
computer this book was checked on, the computer starts at once; Esc as
it starts shows the menu. That is the other system's GRUB setting
(\texttt{GRUB\_TIMEOUT\_STYLE=hidden}, \texttt{GRUB\_SAVEDEFAULT=true}
in \pth{/etc/default/grub}), and the installer leaves it as it finds it.

The whole system is read into RAM at boot here too --- everything on the
\texttt{K4510} partition, which is why nothing large should be kept on it
beside the system itself.

\section{A Tektronix beside the machine}
\texttt{tek40xx} (Ian Schofield's, GPL-3) is a Tektronix 4010/4014
storage-tube terminal that is also a telnet client. It is on the K4510's
Linux and in the container because of what a storage tube is
\emph{for}: a PiDP-11 answers on telnet ports, and a machine of that
vintage deserves a screen of that vintage to draw on.

\begin{type}
tek HOST [PORT]
\end{type}

\noindent on the second console --- full screen on a bare console, in a
window under a desktop. And without any PDP-11 at all:

\begin{type}
tekplay NAME|FILE...
\end{type}

\noindent plays Tektronix plot files at a 9600-baud pace, which is the
right speed to watch one being drawn. With no argument it plays every
plot that ships; \texttt{HOME} clears the screen and \texttt{END} quits.
Nothing in the emulator changed for any of this: it is a second program
on the same computer, which is exactly what a second terminal was in
1975.

\section{From another computer}
The machine's Linux answers ssh (as the user \texttt{k4510}), which is
how the menu file is edited once the machine is locked, and how a
screenshot is taken from across the room:

\begin{type}
ssh k4510@machine k4510-shot
\end{type}

\noindent writes one, as PrtSc would, into \pth{shots/}. And the
machine can be typed into from there, as if at its keyboard:

\begin{type}
ssh k4510@machine "k4510-type 'BOOK\n'"
ssh k4510@machine k4510-type --key down down enter
\end{type}

\noindent \verb|\n| is Enter and \verb|~| a half-second pause; keys
without a letter go by name (\verb|up|, \verb|pgdn|, \verb|esc|,
\verb|f1|\ldots). Typing and a screenshot at a time is the machine
driven from another room --- F12's menu included, with
\verb|--key f12|. F12 $\rightarrow$ Input $\rightarrow$ \emph{Key pipe}
can turn it off, or --- as it starts --- show every key so typed at the
foot of the screen. \verb|k4510-screen| is the third tool beside them:
the machine's text screen, as text, so another computer can read what
the machine says without looking at a picture.

On the computer doing the driving, \pth{tools/k4510-remote} puts the
three together behind one command --- \verb|status|, \verb|type|,
\verb|key|, \verb|screen|, \verb|expect| (wait until a text appears),
\verb|shot|, \verb|run| (a script of those, one a line), \verb|logs| and
\verb|reboot| --- and \verb|k4510-remote tui| is all of it from a menu,
with the machine's screen live above it and a pass-through mode in which
what you type goes straight to the machine (Ctrl-] leaves).
